\section{基本函数}
\label{sec:ru-elementary-functions}

后面的例子反复用到几类具体函数：多项式与有理函数、幂函数、指数与对数、三角函数及其反函数。
它们在中学已经出现过，本节只复习怎样辨认它们、它们允许哪些输入，以及常用的恒等式；严格的构造与分析性质留给后两章。

形如$p(x)=c_0+c_1x+\cdots+c_nx^n$的函数称为\term{多项式}（polynomial），它对每个实数$x$都有定义；
两个多项式之商$p(x)/q(x)$称为有理函数，定义域要去掉$q$的零点。幂函数$x^\alpha$的允许输入取决于指数：
正整数次幂对所有$x$有定义；负整数次幂如$x^{-1}=1/x$要求$x\neq0$；$x^{1/2}=\sqrt x$只对$x\geq0$取实数值；
指数$\alpha$不是有理数时，$x^\alpha$只对$x>0$定义，第~\ref{sec:ru-elementary-derivatives}节将用指数和对数给出它的定义。
所以同一个记号$x^\alpha$，定义域会随$\alpha$改变。

设$a>0$且$a\neq1$。指数函数$a^x$对所有实数$x$有定义，取值恒为正，满足$a^{x+y}=a^xa^y$、$(a^x)^y=a^{xy}$；
$a>1$时它严格递增，$0<a<1$时严格递减。由上一节，严格单调的函数有反函数，指数函数的反函数就是以$a$为底的对数$\log_ax$，定义域为$(0,+\infty)$。
两者互为反函数，意味着$a^{\log_ax}=x$（$x>0$）与$\log_a(a^x)=x$（$x\in\mathbb R$），由此得到$\log_a(xy)=\log_ax+\log_ay$。
底数取$e\approx2.71828$时，$e^x$称为自然指数函数，其反函数记作$\log x$（也记$\ln x$），称为自然对数；本书不加底数的$\log$都指自然对数。
底数为何要求$a>0$：负数的分数次幂可能不是实数，例如$(-1)^{1/2}$；$a=1$时$1^x\equiv1$不是单射，没有反函数。

三角函数通过单位圆定义。从点$(1,0)$出发，沿单位圆逆时针走过弧长$\theta$，到达的点记作$(\cos\theta,\sin\theta)$；
$\theta<0$时顺时针走。用弧长度量角称为弧度制，一周的弧长为$2\pi$。由这一定义，$\cos^2\theta+\sin^2\theta=1$，
正弦和余弦以$2\pi$为周期，正弦是奇函数、余弦是偶函数。正切$\tan\theta=\sin\theta/\cos\theta$要求$\cos\theta\neq0$，即$\theta\neq\pi/2+k\pi$，它以$\pi$为周期。
加法公式
\[
 \sin(\alpha+\beta)=\sin\alpha\cos\beta+\cos\alpha\sin\beta,\qquad
 \cos(\alpha+\beta)=\cos\alpha\cos\beta-\sin\alpha\sin\beta
\]
在复数乘法中已经用过，后面求三角函数的导数时还要用到。

三角函数是周期函数，在整个定义域上不是单射。要得到反函数，须先选一个函数严格单调、并取遍全部函数值的区间。
正弦在$[-\pi/2,\pi/2]$上从$-1$严格增到$1$，它在这一区间上的反函数记作$\arcsin:[-1,1]\to[-\pi/2,\pi/2]$；
余弦在$[0,\pi]$上从$1$严格减到$-1$，相应的反函数为$\arccos:[-1,1]\to[0,\pi]$；正切在$(-\pi/2,\pi/2)$上严格递增，反函数为$\arctan:\mathbb R\to(-\pi/2,\pi/2)$。
例如$\sin(5\pi/6)=1/2$，但$\arcsin(1/2)=\pi/6$，因为$5\pi/6$不在选定的区间内。表~\ref{tab:ru-elementary-functions}汇总这些函数的定义域、值域与周期。

\begin{table}[htbp]
 \centering\small
 \begin{tabular}{@{}llll@{}}
 \toprule
 函数 & 定义域 & 值域 & 周期\\
 \midrule
 $x^n$（$n$为正奇数/偶数） & $\mathbb R$ & $\mathbb R$\,/\,$[0,+\infty)$ & \\
 $x^{-n}$（$n$为正奇数/偶数） & $x\neq0$ & $x\neq0$\,/\,$(0,+\infty)$ & \\
 $\sqrt x$ & $[0,+\infty)$ & $[0,+\infty)$ & \\
 $a^x$（$a>0$，$a\neq1$） & $\mathbb R$ & $(0,+\infty)$ & \\
 $\log_ax$（$a>0$，$a\neq1$） & $(0,+\infty)$ & $\mathbb R$ & \\
 $\sin x$，$\cos x$ & $\mathbb R$ & $[-1,1]$ & $2\pi$\\
 $\tan x$ & $x\neq\pi/2+k\pi$ & $\mathbb R$ & $\pi$\\
 $\arcsin x$ & $[-1,1]$ & $[-\pi/2,\pi/2]$ & \\
 $\arccos x$ & $[-1,1]$ & $[0,\pi]$ & \\
 $\arctan x$ & $\mathbb R$ & $(-\pi/2,\pi/2)$ & \\
 \bottomrule
 \end{tabular}
 \caption{基本函数的定义域、值域与周期。斜线前后分别对应奇数与偶数次幂；$\sqrt x$是平方函数在$[0,+\infty)$上的反函数，$a^x$与$\log_ax$互为反函数。三角函数采用弧度制，$\log$不写底数时指自然对数。}
 \label{tab:ru-elementary-functions}
\end{table}

需要说明这些描述依赖什么。上面把指数函数当作已知对象使用，但“$a^x$对无理数$x$是什么数”“$e$是什么数”都需要极限才能回答；
单位圆上的弧长、$\sin h$与$h$在$h$很小时的比较，也需要极限。
下一章将先用级数构造自然指数函数，由代数和估计得到它的正性、单调性与加法公式，
在介值定理之后严格构造自然对数，并证明三角函数的基本极限；第~\ref{sec:ru-euler-formula}节再由级数得到Euler公式。
在那之前，本节的内容只作为辨认和计算的背景，不用来论证无限求和、连续性或指数与三角函数的关系。
